\section{The mathematical definition}\label{sec:06-groups:01-definition:1}

The checkpoint project of Chapter 5 asked for a \texttt{Monoid}, a set with an associative operation and a two-sided identity, built entirely from the bundling techniques of Chapters 1--5, with no inverses required. A group is the natural next question a monoid forces: what happens once every element is also required to be \emph{undoable}? Adding that one further axiom is the entire content of the opening definition of this chapter, and it is what turns ``a set with an associative operation'' into the richer structure, with cancellation, and eventually genuine theorems reusable across every group at once, that the rest of Part II builds on.

A \textbf{group} is a set \(G\) together with:

\begin{itemize}
\tightlist
\item
  a binary operation \(\cdot : G \times G \to G\),
\item
  a distinguished element \(e \in G\) (the identity),
\item
  an inverse function \((-)^{-1} : G \to G\),
\end{itemize}

satisfying three axioms, for all \(a, b, c \in G\):

\[
\begin{aligned}
\text{(associativity)} &\quad (a \cdot b) \cdot c = a \cdot (b \cdot c) \\
\text{(identity)} &\quad e \cdot a = a \quad\text{and}\quad a \cdot e = a \\
\text{(inverse)} &\quad a^{-1} \cdot a = e \quad\text{and}\quad a \cdot a^{-1} = e
\end{aligned}
\]

No further axiom is required; in particular, commutativity is not assumed in general (a group where \(a \cdot b = b \cdot a\) always holds is called \textbf{abelian} or \textbf{commutative}, defined separately below).

\subsection{Sources, quoted}\label{sec:06-groups:01-definition:2}

Formal definitions and citations for this section, gathered here for reference (full entries in the \hyperref[sec:root:bibliography:1]{Bibliography}):

\begin{itemize}
\tightlist
\item
  \textbf{Group.} A set \(G\) with a binary operation, identity, and inverses satisfying associativity and the identity/inverse laws (below); consequences of these axioms, e.g.~``the identity of \(G\) is unique \ldots{} for each \(a \in G\), \(a^{-1}\) is uniquely determined'' (\cite{DummitFoote2003}, §1.1 ``Basic Axioms and Examples,'' p.~17, Proposition 1).
\item
  Aluffi (\cite{Aluffi2009}) is offered as further reading, not an independently verified factual claim. The use of forgetful functors and universal properties by Aluffi is publicly documented in the table of contents of the book itself, not quoted from a verified excerpt.
\end{itemize}
